% Language: LaTeX
% =====================================================================================================
%  formato.tex  --  Configuracion de apariencia del documento
%  (margenes, espaciado, titulos, encabezados, indices)
% =====================================================================================================

% --- Margenes ---
\geometry{a4paper, left=2.5cm, right=2.5cm, top=3.5cm, bottom=3.5cm}

% --- Interlineado ---
\onehalfspacing

% --- Profundidad de numeracion / indice (subsubsections) ---
\setcounter{tocdepth}{4}
\setcounter{secnumdepth}{4}
\counterwithin{table}{section}
\counterwithin{figure}{section}

% --- Titulo (caption) ARRIBA en tablas ---
% 'plaintop' coloca el caption sobre el cuerpo del flotante sin importar
% donde este escrito el \caption en el fuente. Requiere 'float' (paquetes.tex).
\floatstyle{plaintop}
\restylefloat{table}
\floatstyle{plain}                 % restaura el estilo por defecto para lo demas
% \captionsetup[table]{position=top, skip=6pt}
\captionsetup[table]{position=top, skip=6pt, labelsep=period}
\captionsetup[figure]{labelsep=period}

% --- Titulos: SECTION ---
\titleformat{\section}
  {\normalfont\bfseries\large}{\thesection}{1em}{\MakeUppercase}
\titlespacing*{\section}{0pt}{6pt}{18pt}
\newcommand{\sectionbreak}{\clearpage}   % cada seccion en pagina nueva

% --- Titulos: SUBSECTION ---
\titleformat{\subsection}
  {\normalfont\bfseries\normalsize}{\thesubsection}{1em}{}
\titlespacing*{\subsection}{0pt}{6pt}{18pt}

% --- Titulos: SUBSUBSECTION ---
\titleformat{\subsubsection}
  {\normalfont\bfseries\normalsize}{\thesubsubsection}{1em}{}
\titlespacing*{\subsubsection}{0pt}{6pt}{18pt}

% --- Titulos: PARAGRAPH (4to nivel, actua como subsubsubsection) ---
% titleformat con forma [hang] (default) lo convierte de run-in a titulo de bloque
\titleformat{\paragraph}
  {\normalfont\bfseries\normalsize}{\theparagraph}{1em}{}
\titlespacing*{\paragraph}{0pt}{6pt}{18pt}
% --- Alias: nomenclatura consistente con la jerarquia de secciones ---
\newcommand{\subsubsubsection}{\paragraph}

% --- Resumen / Abstract (seccion sin numero, centrada) ---
\titleformat{name=\section,numberless}[block]
  {\normalfont\bfseries\normalsize\centering}{}{0pt}{}

% --- Renombrar indices en espanol ---
% \addto\captionsspanish{
%   \renewcommand{\contentsname}{Indice General}
%   \renewcommand{\listfigurename}{Indice de Imagenes}
%   \renewcommand{\listtablename}{Indice de Figuras}
%   \renewcommand{\tablename}{Tabla}
%   \renewcommand{\refname}{Referencias Bibliograficas}
% }

\addto\captionsspanish{
  \renewcommand{\contentsname}{Índice General}
  \renewcommand{\listtablename}{Índice de Tablas}
  \renewcommand{\listfigurename}{Índice de Figuras}
  \renewcommand{\tablename}{Tabla}
  \renewcommand{\figurename}{Figura}
  \renewcommand{\refname}{Referencias Bibliográficas}
}

% --- Encabezados y pies (fancyhdr) ---
\pagestyle{fancy}
\fancyhf{}
\fancyhead{}
\fancyfoot[R]{\thepage}
\renewcommand{\headrulewidth}{0pt}
\renewcommand{\footrulewidth}{0pt}

% Paginas preliminares en numeracion romana minuscula
\fancypagestyle{plain}{
  \fancyhf{}
  \fancyfoot[R]{\thepage}
  \renewcommand{\thepage}{\roman{page}}
}
% Portada de grupo para Apendices / Anexos (titulo VISIBLE + entrada en indice)
\newcommand{\grupoapendices}[1]{%
  \clearpage              % empieza en pagina nueva
  \phantomsection         % ancla correcta para el hipervinculo del TOC (hyperref)
  \section*{#1}%          % titulo VISIBLE en la pagina (no numerado)
  \addcontentsline{toc}{section}{#1}%   % entrada en el indice
}

% --- Figura a pagina completa (recupera margen superior/inferior) ---
% Uso: \begin{figuracompleta}[top=1cm, bottom=1.5cm] ... \end{figuracompleta}
% Ocupa una pagina propia; los margenes se restauran al salir.
% \newenvironment{figuracompleta}[1][top=1.5cm, bottom=1.5cm]{%
%   \clearpage
%   \newgeometry{left=2.5cm, right=2.5cm, #1}%
%   \thispagestyle{fancy}%  numero de pagina visible; usar empty para quitarlo
%   \noindent
%   \begin{minipage}{\linewidth}
%   \centering
% }{%
%   \end{minipage}
%   \clearpage
%   \restoregeometry
% }



% \newenvironment{figuracompleta}[1][top=1.5cm, bottom=1.5cm]{%
%   \clearpage
%   \newgeometry{left=2.5cm, right=2.5cm, #1}%
%   \thispagestyle{fancy}%  numero de pagina visible; usar empty para quitarlo
%   \noindent
%   \begin{minipage}{\linewidth}
%   \captionsetup{type=figure}%
%   \centering
% }{%
%   \end{minipage}
%   \clearpage
%   \restoregeometry
% }

\newenvironment{figuracompleta}[1][top=1.5cm, bottom=1.5cm]{%
  \clearpage
  \newgeometry{left=2.5cm, right=2.5cm, #1}%
  \thispagestyle{fancy}%  numero de pagina visible; usar empty para quitarlo
  \noindent
  \begin{minipage}{\linewidth}
  \captionsetup{type=figure}%
  \centering
}{%
  \end{minipage}
  \clearpage
  \aftergroup\restoregeometry
}